\documentclass{article}
\usepackage[T1]{fontenc}
\usepackage{lmodern}
\usepackage{graphicx}
\usepackage{amsmath,amssymb}
\usepackage[a4paper,margin=2cm]{geometry}
\usepackage[absolute,overlay]{textpos}
\setlength{\TPHorizModule}{1mm}
\setlength{\TPVertModule}{1mm}
\usepackage[indonesian]{babel}
\usepackage{hyperref}
\setlength{\emergencystretch}{2em}
% unit: Halaman 9

\begin{document}

% === Halaman 9 (llm) ===

\section*{Enkripsi}

\begin{enumerate}
    \item \textit{Preprocessing}: Pesan yang akan dienkripsi dikodekan menjadi integer $m$. Jika $m$ berukuran besar, pecah $m$ menjadi blok-blok plainteks yang lebih kecil: $m_1, m_2, m_3, \dots$ (syarat: $0 \le m_i < n - 1$)
    \item Hitung cipherteks $c$ untuk plainteks $m$ menggunakan kunci publik $e$ dengan persamaan
    \[ c = m^e \text{ mod } n \]
    Jika $m$ dipecah menjadi $m_1, m_2, m_3, \dots$, enkripsi setiap $m_i$ menjadi
    \[ c_i = m_i^e \text{ mod } n \]
\end{enumerate}

\end{document}
